\section{Introduction}

An Indian commercial electricity bill has three levers: a time-of-day energy
charge, a demand charge levied on the highest 30-minute average demand in the
month, and power-factor penalties. The second dominates the risk. It is set by
one block out of roughly $1{,}440$ in a month, so a single bad afternoon prices
all thirty days. A controller that reduces average consumption but occasionally
overshoots can cost more than one that does nothing.

The standard applied response is a chance constraint: replace the forecast mean
in the capacity constraint with a high quantile, typically the 95th, and treat
the result as a $95\%$ guarantee. Formally, with $b_t$ the uncontrollable base
load, $u_t$ the controllable allocation, $s_t$ solar generation and $D$ the
ceiling,
%
\begin{equation}
  \hat b^{\,q95}_t + u_t - \hat s^{\,q05}_t \;\le\; D
  \qquad \text{for each } t \text{ in the horizon.}
  \label{eq:marginal}
\end{equation}
%
Equation~\eqref{eq:marginal} is a set of \emph{marginal} statements. Each row
says that at that timestep, considered alone, the realisation exceeds the bound
about $5\%$ of the time. The event that costs money is different: it is
%
\begin{equation}
  \Pr\Big[\max_{t \in \mathcal{W}} \big(b_t + u_t - s_t\big) > D\Big],
  \label{eq:joint}
\end{equation}
%
a statement about a maximum over a window $\mathcal{W}$. That
\eqref{eq:marginal} does not imply \eqref{eq:joint} is elementary and well
known; joint chance constraints, their intractability, and the standard routes
around them --- Boole/Bonferroni risk allocation and the scenario approach
\citep{calafiore2006scenario,nemirovski2006convex} --- are textbook material.

\paragraph{What is not established} is how large the gap is on real building
load, and how badly each of the standard substitutes misses it. Both ends matter
practically: the first is how much a deployed controller is under-protected, and
the second is how much headroom a careful designer gives away for nothing. We
measure both, and find they are wrong in opposite directions by factors we can
state.

\paragraph{Contributions.}
\begin{enumerate}
  \item A measured bracket on horizon risk for a demand ceiling on held-out
        meter data (Section~\ref{sec:horizon}), over fourteen windows in two
        countries: claimed $0.05$, realised $0.39$--$0.82$ (median $0.57$;
        no $95\%$ lower bound under $0.28$), independence calculation
        $0.87$--$0.99$, union bound vacuous on every row, and the gap surviving
        a restatement of the per-step layer's own tuning parameter in units
        that transfer across meter sizes. The common formulation is under-protected by a factor of eight to
        sixteen, the usual correction is $1.7\times$ too conservative at the
        median, and the gap is a distribution rather than a number.
  \item A closed-loop demonstration, on two control objects in two countries
        under their own tariffs, that this risk level can be made an
        operator-set parameter rather than an emergent property of horizon
        length --- with the half of that claim which replicates separated from
        the half which does not. The response is monotone on both (rank
        correlation $0.957$ and $0.988$); it is close to exact and errs safe
        only on the plant whose per-step layer is itself calibrated
        (Section~\ref{sec:dial}).
  \item A finding made and then unmade by its own audit. One protocol over
        thirty supplies in eight countries, replicated inside a single city
        across building, campus and system demand, showed conformal coverage
        degrading with aggregation. The audit of the adaptive layer at system
        scale showed why: its step is stated in the units of the series and had
        stopped moving. Re-read without that confound, split conformal fails at
        every tier alike and the adaptive layer, with its step as a fraction of
        the interval width, holds the level at every tier
        (Section~\ref{sec:calibration}). What transfers across scales is not
        the number but the fraction, and a practitioner moving the layer from a
        building to a feeder will hit exactly this.
  \item A null separating forecast skill from available flexibility as
        independent deployment axes, replicated on a panel it was not fitted to
        (Section~\ref{sec:axes}).
  \item A provenance audit of a dataset added to the study under an external
        constraint, which disqualifies that arm's most favourable result and is
        reported as such (Section~\ref{sec:china}).
\end{enumerate}

\paragraph{Where the novelty is, and where it is not.} We invented no method.
Every component --- chance-constrained MPC for buildings, joint chance
constraints and their standard treatments, quantile gradient boosting,
conformalized quantile regression, adaptive conformal inference, copula
scenario generation, decision-focused learning --- is prior work, and
Section~\ref{sec:related} places each. The novelty is empirical and
integrative: an instrument that pins everything and moves one thing, run on
metered load, with the results that argue against it reported at full size.
Section~\ref{sec:related} states the gap that survives the literature.
